% Part: normal-modal-logic
% Chapter: axioms-systems
% Section: duals

\documentclass[../../../include/open-logic-section]{subfiles}

\begin{document}

\olfileid{nml}{prf}{dua}

\olsection{ద్వంద్వ \usetoken{P}{formula}}

\begin{defn}\ollabel{def:duals}
  \Ax{T}, \Ax{B}, \Ax{4}, \Ax{5} అనే \tetoken{సూత్రాలలో}{!!{formula}s}
  ఒక్కొక్కదానికి ఒక \emph{ద్వంద్వ సూత్రం} ఉంది. దాన్ని ఉపసూచికలో
  డైమండ్ గుర్తుతో ఈ విధంగా సూచిస్తాం:
  \begin{align}
    \tag{\Ax{T_\Diamond}} p & \lif \Diamond p\\
    \tag{\Ax{B_\Diamond}} \Diamond\Box p & \lif p\\
    \tag{\Ax{4_\Diamond}} \Diamond\Diamond p & \lif \Diamond p\\
    \tag{\Ax{5_\Diamond}} \Diamond\Box p & \lif \Box p
    \end{align}
\end{defn}

పై ద్వంద్వ \tetoken{సూత్రాలలో}{!!{formula}s} ఒక్కొక్కటి, దానికి
సంబంధించిన \tetoken{సూత్రంలో}{!!{formula}} $p$ స్థానంలో $\lnot p$ను
ప్రతిస్థాపించి, ప్రతిలోమం తీసుకుని, $\lnot\Box\lnot$ స్థానంలో
$\Diamond$ను, $\lnot\Diamond\lnot$ స్థానంలో~$\Box$ను ఉంచడం ద్వారా
లభిస్తుంది. \Ax{D}, అంటే $\Box!A \lif \Diamond!A$, ఈ అర్థంలో
తనకుతానే ద్వంద్వం.

\begin{prop}\ollabel{prop:dualsys}
  \olref[nml][prf][dua]{def:duals}లోని ప్రతి
  \tetoken{సూత్రం}{!!{formula}}~$!A$కు:
  $\Log{K}!A = \Log{K}!A_{\Diamond}$.
\end{prop}
\begin{proof}
  వ్యాయామం.
\end{proof}
\begin{prob}
  \olref[nml][prf][dua]{prop:dualsys}ను నిరూపించండి.
\end{prob}

\end{document}
